\section{The Eisenstein proof}\label{sec:eisenstein}
\subsection{Arithmetic preliminaries and split-prime cores}
Let $Q(x,y)=x^2-xy+y^2$, $\omega^2+\omega+1=0$ and
$R=\Z[\omega]$.  Thus $Q(x,y)=N(x+y\omega)$.  We use the
semigroup $\AA$ defined above and write $\NN=P_{-3}$.
Every member of $\AA$ is $1\bmod6$.  The ramified dilation
$3\!\cdot\!\AA=\{3a:a\in\AA\}$ has residue $3\bmod6$.

The elementary norm-value identity separates the two binary problems.
Its proof uses only prime splitting and unique factorization.

\begin{proposition}\label{prop:primitive-classification}
The primitive Eisenstein norms are exactly
\[
\NN=\AA\sqcup(3\!\cdot\!\AA).
\]
\end{proposition}
\begin{proof}
If an inert prime $p$ divides $Q(x,y)$, it divides both coordinates,
since $R/pR$ is a field, including when $p=2$. Also $9\mid Q(x,y)$ forces $3\mid x,y$.
Indeed, $3\mid Q(x,y)$ gives $x+y\equiv0\pmod3$; writing
$y=-x+3t$ gives
\[
Q(x,y)=3(x^2-3xt+3t^2).
\]
Divisibility by $9$ now forces $3\mid x$, and hence $3\mid y$.
Thus a primitive norm has no inert prime factor and contains at most
one factor $3$. It belongs to the stated union.

Conversely, the triangular Eisenstein lattice has covering radius
$1/\sqrt3$, so rounding gives a remainder of norm at most $1/3<1$.
Thus $R=\Z[\omega]$ is norm Euclidean and is a unique factorization
domain. For a rational
prime $q\equiv1\pmod3$, choose a nontrivial cube root
$t\in\F_q^\times$. The ideal $(q,\omega-t)$ has quotient $\F_q$;
its principal generator $\pi_q$ therefore has norm $q$. The two distinct
roots $t,t^{-1}$ yield conjugate, nonassociate factors. Thus
$q=\pi_q\bar\pi_q$, up to a unit.
For $a=\prod q^{e_q}\in\AA$, choose one factor $\pi_q$ at each $q$ and
put $z=\prod\pi_q^{e_q}$. Then $N(z)=a$. The element
$(1-\omega)z$ has norm $3a$. Neither element is divisible by a rational
prime: a split rational prime would require both conjugate factors,
an inert prime cannot divide its norm, and $3$ is associated to
$(1-\omega)^2$, whereas at most one factor $1-\omega$ was used.
Their coordinates in the integral basis $1,\omega$ are therefore
coprime. Finally, the two sets are disjoint, since elements of $\AA$
are prime to $3$.
\end{proof}

\begin{theorem}\label{thm:main}
One has
\[
 \AA+\AA=\{n>0:n\equiv2\pmod6\},\qquad
 \AA+(3\!\cdot\!\AA)=\{n>0:n\equiv4\pmod6\}.
\]
\end{theorem}

\paragraph{A reduction that helps both problems.}
Put $\PP=\{p\text{ prime}:p\equiv2\pmod3\}$. For either admissible residue
class, write $n=gn_0$, where $g\in\AA$ contains every prime-power factor
of $n$ with prime $1$ modulo $3$. Since $g\equiv1\pmod6$, the target
$n_0$ has the same residue as $n$. A decomposition of either type for
$n_0$ scales by $g$ to the corresponding decomposition of $n$.
We can therefore concentrate on
\begin{equation}\label{eq:reduced}
n\equiv2\text{ or }4\pmod6,\qquad
p\mid n\Longrightarrow p\in\PP.
\end{equation}
This reduction will matter in the divisor estimates and in the exact
calculation at $3$. The finite certificates cover all targets in their
respective ranges, so they need no reduction.
